\documentclass[15pt]{article}
\usepackage[utf8]{inputenc}

\usepackage{geometry}
\geometry{left=1.2in, right=1.2in, top=.72in, bottom=.72in}
\usepackage{amsmath}
\usepackage{amssymb}
\usepackage[colorlinks=true]{hyperref}
\usepackage{titlesec}
\titleformat*{\section}{\bfseries\large}
\newtheorem{thm}{Theorem}
\newcommand{\xn}{(x_n)_{n \in \mathbb{N}} }



\title{Genova - Graz Number Theory Workshop\\ {\large May 30 - June 1, 2022}\\[1em] Abstracts}
\author{}
\date{}

\uchyph=0

\begin{document}

\maketitle
\thispagestyle{empty}

\vspace{-5em}

\subsection*{Bence Borda: Randomness of Sudler products}

We discuss recent results and conjectures on the random behavior of the Sudler product $P_N(\alpha)=\prod_{n=1}^N |2 \sin (\pi n \alpha)|$. Complementing a theorem of J. Beck, we show that $\log P_N(\alpha)$ satisfies a so-called temporal central limit theorem at certain quadratic irrationals as N ranges in a long interval of integers, and extend this to more general Birkhoff sums for the irrational rotation. Sudler products are also related to the Kashaev invariant of the figure-eight knot in algebraic topology, which satisfies a remarkable asymptotic modularity relation. A recent solution of a conjecture of Zagier on the Kashaev invariant has led to its limit distribution at typical rational arguments, as conjectured by Bettin and Drappeau.



\subsection*{Emanuel Carneiro: Some Fourier optimization problems in number theory}

This will be a talk on how some Fourier optimization problems appear naturally in connection to problems in number theory. We will discuss recent new bounds related to prime gaps, to the pair correlation conjecture, and to low-lying zeros of families of $L$-functions.

\subsection*{Sneha Chaubey: Zeros of the Derivatives of the Selberg class L functions}

Inspired from Spieser's and Spira's results on the relation between zeros of the Riemann zeta function and it's derivatives, in 1974, Levinson and Montgomery came up with several results concerning the zeros of the $k$-th derivatives of the Riemann zeta.  In one such result, they show that the RH implies that all the $k$-th derivatives of the Riemann Zeta function have at most a finite number of non-real zeros in the left of the critical line. In this talk, we discuss this for L functions belonging to the Selberg class and prove that an analogous result holds true for this wide class of functions as well. This is a joint work with Suraj Singh Khurana and Ade Irma Suriajaya.

\subsection*{Sary Drappeau: Quantum modularity for central values of additively twisted Maass $L$-functions}

Let $d(n)$ be the divisor function and for all rational $x$, define
$$D(x, s) = \sum_{n\geq 1}  d(n) e^{2\pi i n x} / n^s \qquad  (\Re(s)>1).$$
This can be seen as the linear additive twist of  $\zeta^2$.
For rational $x$, we know that $D(x, s)$ extends to a meromorphic function of $s$, regular at $1/2$.
In 2016, Bettin proved that the central value $D(x) := D(x, 1/2)$, as a function of $x$, is a ``quantum modular" as defined by Zagier, in the precise sense that the map
$$h(x) = D(x) - D(-1/x)$$
initially defined on ${\mathbb Q}^*$ extends to a H\"older-continuous map on ${\mathbb R}^*$.
A similar property holds true for modular symbols attached to holomorphic cusp forms.
These properties are of interest when looking at the value distribution of these maps over the set of rationals with bounded denominators, a question studied in various aspects with various tools by Petridis-Risager, Nordentoft, Lee-Sun and Bettin-D..

In this talk we discuss a generalization of the above properties for Maass forms on Fuchsian groups, which conjecturally includes all cases of degree 2 $L$-functions. We discuss related applications.
This is joint work with A. Nordentoft.


\subsection*{Daniel El-Baz: Primitive rational points on expanding horospheres: effective joint equidistribution}

I will report on ongoing work with Min Lee and Andreas
Str\"ombergsson. Using techniques from analytic number theory, spectral
theory, geometry of numbers as well as a healthy dose of linear algebra
and building on a previous work by Bingrong Huang, Min Lee and myself,
we furnish a new proof of a 2016 theorem by Einsiedler, Mozes, Shah and
Shapira. That theorem concerns the equidistribution of primitive
rational points on certain manifolds and our proof has the added benefit
of yielding a rate of convergence. It turns out to have (perhaps
surprising) applications to the theory of random graphs, which I shall
also discuss.


\subsection*{Jungwon Lee: Another view of Ferrero-Washington Theorem}

We reprove the main equidistribution result in the Ferrero--Washington proof of the vanishing of cyclotomic Iwasawa $\mu$-invariant, based on the ergodicity of a certain $p$-adic skew-extension dynamical system that can be identified with Bernoulli shift (joint with Bharathwaj Palvannan).  

\subsection*{Marc Munsch: Zeros of Fekete polynomials and applications}

The study of the location of zeros of polynomials with coefficients 
constrained in different sets has a very rich history. Fekete 
polynomials are certain Littlewood type polynomials whose coefficients 
are the values of a quadratic Dirichlet character. These polynomials 
were considered by Fekete in an attempt to understand real zeros of 
Dirichlet L-functions in the critical strip. Since then, their zeros 
and value distribution have been extensively studied.

I will present some recent results concerning zeros of Fekete 
polynomials and discuss applications to sign changes of quadratic 
character sums. This is a joint work with O. Klurman and Y. Lamzouri.


\subsection*{Marc Technau: Averages of Dedekind sums and connections to Euclid-type algorithms and continued fractions}
Investigation into averages of Dedekind sums has attracted some interest. For instance, Conrey, Fransen, Klein and Scott (1996) have obtained asymptotic formulae for $\ell^p$ averages of Dedekind sums over Farey points with even p. Girstmair and Schoi\ss engeier (2005) have determined the correct order of magnitude of the same quantity for p=1, yet their method falls short of establishing a full asymptotic formula.

In this talk, we aim to survey some works around these and related topics. Moreover, we intend to discuss recent joint work with Athanasios Sourmelidis and Paolo Minelli on restricted averages of Dedekind sums. This gives a quantitative answer to a irregularities of distribution of Dedekind sums observed by H. Ito (2004) on the basis of numerical computations. By classical connections between Dedekind sums, a Euclid-type algorithm and continued fraction expansions, the above has a nice interpretation concerning irregularities of distribution of average running times and certain averages of partial quotients. The method of proof is based on the above connections and averages of lengths of continued fraction expansions, which is well-studied topic in the literature. A similar approach may also yield further insights into $\ell^1$ averages of Dedekind sums.


\subsection*{Francesco Veneziano: Rational angles between points in a plane lattice}

Given a plane lattice $\Lambda$ we study the triples of elements $A,B,C\in
\Lambda$ such that the angle $\hat{ABC}$  is a rational multiple of $\pi$.
I will present a classification of the lattices and configurations which
arise. This is a joint work with R. Dvornicich, D. Lombardo and U. Zannier

\subsection*{Agamemnon Zafeiropoulos: Poissonian correlations: higher orders and weaker variants}

Let $\xn \subseteq [0,1]$ be sequence of points in the unit interval. We say that $\xn$ has Poissonian pair correlations (PPC) if \[ \lim_{N\to \infty} \frac{1}{N}\#\Big\{ m,n\leqslant N, m\neq n : \|x_m-x_n\| \leqslant \frac{s}{N} \Big\} = 2s  \]
for all $s>0.$ It is known that sequences with PPC are also uniformly distributed. We show that the same conclusion is true for sequences with Poissonian correlations of any order $k\geqslant 3.$ Moreover we define weaker variants of the notion of PPC and examine their relations with equidistribution. (Joint work with M. Hauke.) 

\end{document}
